% Generated draft: requires mathematical and source review.
\documentclass[11pt]{article}
\usepackage[margin=25mm]{geometry}
\usepackage{fontspec}
\setmainfont{Latin Modern Roman}
\usepackage{amsmath,amssymb,mathtools,unicode-math}
\setmathfont{Latin Modern Math}
\usepackage{enumitem,needspace,xurl,hyperref}
\setlength{\parindent}{0pt}
\setlength{\parskip}{5pt}
\setlength{\emergencystretch}{3em}
\newcommand{\N}{\mathbb{N}}
\newcommand{\Z}{\mathbb{Z}}
\newcommand{\Q}{\mathbb{Q}}
\newcommand{\R}{\mathbb{R}}
\newcommand{\C}{\mathbb{C}}
\newcommand{\F}{\mathbb{F}}
\newcommand{\A}{\mathbb{A}}
\newcommand{\PP}{\mathbb{P}}
\newcommand{\E}{\mathbb{E}}
\newcommand{\eps}{\varepsilon}
\newcommand{\dd}{\,\mathrm{d}}
\begin{document}
\section*{UnsolvedMath ID: 30002295}
\textbf{Source problem code:} catalogue entry
\subsection{Definitions and mathematical statement}
\textbf{Research problem.} Derive an analogue of Einstein's formula for the dynamics of localized objects within a fluid-mechanics framework. Concentrating solutions of the Euler equations are suggested as one possible framework.

\textbf{Unresolved formulation.} The supplied record does not identify the particular Einstein formula to be generalized or specify the mathematical meaning of an analogue. Its supplied literature assessment describes the intended target as an effective inertial relation, but gives no equation defining that relation. Neither ``localized object'' nor ``concentrating solution'' is defined. The fluid model, spatial dimension and domain, boundary conditions, solution class, concentration parameter and limiting procedure, and effective dynamical variables are unspecified. Consequently, no definite identity, convergence assertion, or quantified theorem can be stated from the supplied evidence alone. These choices remain part of the formulation of the research problem; the Euler-equation approach is an example, not a required hypothesis.
\subsection{Short English statement}
Seek a fluid-mechanical analogue of Einstein's formula governing localized objects, potentially using concentrating Euler flows. The supplied record leaves the formula, fluid model, localization mechanism, and desired effective relation unspecified, so this remains a research question requiring further formulation.
\subsection{Sources}
\begin{itemize}
\item Supplied local UnsolvedMath registry record: permanent ID \(30002295\), source problem code \texttt{catalogue entry}, title ``Einstein-Type Dynamics for Concentrated Euler Flows,'' in the Oberwolfach Reports Open Problems collection. The record gives proposed year 2013 and identifies no individual proposer.
\item Supplied report citation: \emph{Geophysical Fluid Dynamics} (2013), \emph{Oberwolfach Reports}, pp.\ 521--577, DOI \texttt{10.4171/owr/2013/10}. Source URL and supplied literature-source link, labeled ``Oberwolfach Report 2013'': \url{https://doi.org/10.4171/owr/2013/10}. The background attributes the problem to workshop open-problem discussions documented in Oberwolfach Reports; no particular page or speaker is identified.
\item The supplied record labels the statement ``exact'' and reports that its mathematical content was checked against the official report, with notation, context, and question form normalized without changing the source-level claim. This is supplied verification metadata, not an independent source check performed for this draft.
\item Literature assessment dated \texttt{2026-08-21}: the supplied curation labels the problem open and reports that no general derivation of an Einstein-type effective inertial relation for localized objects in fluid mechanics was identified. The listed source checked is the report linked above. This is supplied dated curation, not an independently verified literature assessment; the record explicitly requests independent verification of time-sensitive status and cited claims.
\end{itemize}
\end{document}
